\subsection{Implementation and verification workflow}
\label{sec:implementation}
The software separates dispatch-model export from the reference implementation of the
residual-budget selector. The C++ generator implements the policies of
\cref{tab:policies}, preserves named rows and columns, and writes LP text without invoking
an optimizer. The Python module \path{scripts/research/budget_scaling.py} implements
\cref{prop:budget,prop:dyadic} and supplies the budget-aware factors to the enumerated
experiment. The application-budget audit calls the same selector on stored dispatch rows and
projects each final-policy proposal onto its admissible interval. It retains per-row
decisions and checks paired dispatch incumbents in physical units. When no factor changes,
the budgeted export is identical and needs no separate optimization call. The stress
experiment and dispatch audit assess complementary components of the workflow.

\subsubsection{Row-factor selection}
For a finite, nonconstant row, the selector receives its coefficients, a declared budget
$\delta$, an exported-residual hypothesis $\epsilon$, a coefficient window $[\ell,U]$,
and an optional rounding allowance $E$. The dyadic selector also receives the RHS. A GM
candidate minimizes the logarithmic coefficient objective in exact arithmetic; an L2 candidate can be projected
onto the same interval without acquiring that minimax characterization.
The following procedure makes the implementation choices explicit.
\begin{enumerate}[label=\arabic*.,leftmargin=*,itemsep=3pt]
\item Validate finite inputs, positive budget and residual hypothesis, and a nonempty
positive coefficient window. Constant rows require a separate feasibility check.
Compute the nonzero coefficient extremes $m$ and $M$. If $\delta\le E$, return an
abstention with reason \texttt{rounding\_budget\_exhausted}.
\item Form $L=\max\{\epsilon/(\delta-E),\ell/m\}$ and $H=U/M$ as exact rational
expressions in the supplied binary64 inputs. If $L>H$, report incompatibility.
This step tests the specification before attempting to center the coefficients.
\item For an arbitrary binary64 factor, round $L$ upward and $H$ downward. Abstain if
the resulting interval contains no positive finite factor. Evaluate the GM or L2
candidate using logarithms and clip it to that interval. For exact power-of-two export,
instead intersect with the integer exponent bounds of \cref{prop:dyadic} and choose
the admissible exponent minimizing the coefficient objective, comparing candidates
as exact rational numbers.
\item Return the factor, admissible endpoints and a decision reason. Multiply the
whole row and its RHS by that factor, retaining the original row, budget and factor
for verification. In the stress driver, one abstaining row causes the whole model to
abstain before an optimizer is called.
\end{enumerate}

The arbitrary-factor API returns a \texttt{RowDecision} containing the selected factor,
inward-rounded endpoints, unprojected candidate and decision reason. Its \texttt{factor}
field is \texttt{None} on abstention. Distinguishing an empty real
interval from an interval containing no binary64 factor makes the numerical failure
reviewable. Extrema and norm computation require a linear scan of the row; the exact
checks additionally use integer-ratio arithmetic whose cost depends on the operand
bit lengths. No factor-selection step requires an optimization solve.

\subsubsection{Representation and returned-point checks}
Factor admissibility and exact row export are separate properties. Inward rounding
certifies the factor's interval membership for the stored input numbers. With an arbitrary
factor, independently rounded products can still change proportionality. The dyadic
selector checks every coefficient and the RHS against exact rational products and assumes
gradual underflow. The C++ LP serializer uses
\texttt{std::numeric\_limits<double>::max\_digits10} significant digits to preserve each
stored binary64 value on readback. Decimal round-trip preservation avoids an additional
serialization error; it does not undo multiplication rounding.

The budget conclusion remains conditional on the actual exported-row residual.
The stress checker therefore evaluates original-row residuals exactly as rational
expressions in the stored binary64 data and returned coordinates, then checks bounds,
integrality and the rounded assignment against an independently enumerated optimum.
The dispatch checker uses floating-point residuals against retained original rows.
Those checks are useful diagnostics, but they do not certify the intended pre-rounding
physical data or the correctness of a solver's dual bounds.

\begin{table}[tb]
\centering\small
\caption{Software components and their verification responsibilities. Paths are relative to
the accompanying source and data archive.}
\label{tab:software-components}
\begin{tabularx}{\linewidth}{lXX}
\toprule
Component & Implementation & Responsibility\\
\midrule
Dispatch export & C++ \texttt{PreprocesamientoMIP} and \texttt{ProblemaLp}
& Matched scaling policies; whole-row/RHS transformation; decimal serialization\\
Budget selector & \path{scripts/research/budget_scaling.py}
& Admissible factors, abstention reasons, dyadic product checks and exact residuals\\
Stress driver & \path{scripts/research/residual_budget_experiment.py}
& Fixed models, solver settings, returned points and separate verification labels\\
Application audit & \path{scripts/research/oms_dispatch_validation.py}
& Three precision scenarios, factor intervals and physical-unit checks\\
Activation follow-up & \path{exploration/oms-followup-20261014}
& Finite-box export reserves, matched controls and exact feasible witnesses\\
Independent audit & \path{scripts/research/verify_manuscript_claims.py}
& Re-enumeration, saved acceptance labels, native operational logs and matrix comparisons\\
\bottomrule
\end{tabularx}
\end{table}

\subsubsection{Reproducible evaluation}
The artifact retains original inputs, factors or policy configurations, input hashes,
solver versions and settings, and returned points or operational records. Policy comparisons
match coverage, row roles, kernels and safeguards whenever an information effect is tested.
Fixed-basis conditioning is recomputed after a documented solver-style equilibration;
operational evaluation records both completed-run effort and failures on the assigned pool.
The historical numerical-claim audit reads stored data without importing an optimizer or
the experiment driver. The additional dispatch replay checks preserved exports and
incumbents, independently forms rational interval bounds, and recomputes reported
aggregates; its physical-residual replay shares the reference checker. Both require no
optimization calls. Supplementary Material~1 provides the protocols and
commands; Supplementary Material~2 contains the code and analysis inputs.

An isolated C++ integration applies the same selector to the generator's in-memory
rows before export. On six known models it reproduces the Python factors and
binary64 models exactly while reducing standalone preparation cost.
A separate randomized complete-cycle comparison uses the same isolated executable
with the native selector enabled or disabled, matching all ordered model data.
Across 18 model/seed pairs, both implementations return identical points and native
endpoints with identical work; the paired geometric-mean C++/Python full-process
time ratio is 0.824 on the shared host (Supplementary Material~1, Table~S18).
Compilation and prior witness acquisition are separate costs. This measures the
implementation cost of the same scaling policy, with no reduction in solver work.
